\subsection{The landscape, the swampland and the era of precision cosmology --- arXiv:1808.09440}

\textbf{Authors:} Yashar Akrami, Renata Kallosh, Andrei Linde, Valeri Vardanyan

\paragraph{主な主張}
nilpotent多重項を含むde Sitter超重力の構成を擁護し、検討した弦理論起源のquintessence模型には当時の宇宙論観測との強い緊張があると論じる\cite{Akrami:2018ylq}。

\paragraph{新規性}
no-go議論の仮定と超対称性の非線形実現を整理するとともに、指数型ポテンシャルの傾きに対する統計的制約を再評価する。

\paragraph{位置付け}
de Sitter予想への理論的反論と、代替となる暗黒エネルギー模型への観測的検証を結び付ける研究。

\paragraph{コメント（読書メモ）}
量子補正によってanthropic range内の真空が範囲外へ移っても、別の真空が範囲内へ移り得るため、十分に多数の真空を仮定するランドスケープの描像は頑健だ、という著者らの議論に注目した。また、Maldacena--Nu\~nezのno-go定理\cite{Maldacena:2000mw}をde Sitter構成一般の排除に用いることを批判し、KKLTなどの構成がその前提条件を回避していると説明している。元メモの「定理が反証された」は、この前提の回避という意味で整理する。$|\nabla V|/V$の下限$c$について、当時のデータと対象模型では$c>1$が約$3\sigma$で除外される、という解析結果を記録した。
